% chapters/ch02b_defensive_offense_carryback.tex
% Defensive-Offense and Carryback specifications

\section{Defensive-Offense and Carryback Specification}
\label{sec:defensive-offense-carryback}

This sub-document governs how training-context outputs (TTM/PPS,
adversarial role-play, structural exercises) interface with the
operational layer. It is normative: training and operations are
separated, and the only sanctioned re-entry path is via the
Carryback procedure.

\subsection{Definition}
\label{sec:def-offense-definition}

\begin{description}
\item[Operation (Core)] The decision and execution layer. The
``Strong'' face of the system.
\item[Passive Pattern (PP)] Continuous background design pattern;
defensive monitoring only.
\item[Training] The sandbox layer. ``Structure'' and ``defense'' are
analysed here. \emph{Direct application to operations is
prohibited.}
\item[Defensive-Offense] An offensive-pattern recognition and
defensive-response training mode. Operates only on premise of
survival, never as ``attack execution.''
\item[Anti-Contamination] The principle that training must never
contaminate operational judgement.
\end{description}

\subsection{Ruin Premise}
\label{sec:def-offense-ruin}

\begin{description}
\item[Ruin First] Irreversible state preservation is the primary
constraint.
\item[Optionality Preservation] Future selection capacity is
preserved.
\item[Non-Transfer of Agency] Judgement is never transferred to
training-context systems (they are never authoritative).
\item[Anti-Contamination] Training judgement must not leak into
operations.
\end{description}

\subsection{Defensive-Offense Specification}
\label{sec:def-offense-spec}

The Defensive-Offense system is a permitted Passive Pattern and Type
ZERO interface. It is a \emph{recognition-and-defence} layer, not an
attack-execution layer.

\subsubsection{Permitted Defensive-Offense (Allowed)}
\label{sec:def-offense-allowed}

\begin{itemize}
\item Pattern recognition (incoming pressure / coercion /
manipulation).
\item Threat-signal categorisation.
\item Defensive-response design (restricted to Type ZERO categories).
\item Audit logging (record-only).
\item Premise checking (definition / condition / restriction).
\item Reversible selection (slow tempo, reversible expressions).
\end{itemize}

\subsubsection{Disallowed Offense (Prohibited)}
\label{sec:def-offense-prohibited}

\begin{itemize}
\item Active manipulation, persuasion, exploitation.
\item Attack-design, scripted aggressive sequences, dependency
formation.
\item Use of pattern catalogues to engineer compliance.
\item Any training output forwarded to operations without Carryback.
\end{itemize}

\subsection{Architecture (Training Separation)}
\label{sec:def-offense-architecture}

\paragraph{Principle.}
The Core (operations) layer and the Training layer are
\emph{architecturally separated}. They cannot share state or
authority.

\paragraph{Layer roles.}

\begin{table}[ht]
\centering
\small
\begin{tabularx}{\linewidth}{lX}
\toprule
\textbf{Layer} & \textbf{Role} \\
\midrule
Core SOP & State $\to$ permitted action. Permitted output:
shutdown / hold / proposal (with ON guard). \\
Type ZERO & Mode determination and execution gating. \\
Passive Pattern (PP) & Design pattern (defensive monitoring). \\
Training & Structure-only analysis: Defensive-Offense analysis,
defence design, sandbox practice. \\
Carryback & The \emph{only} sanctioned Training $\to$ Core
direction. \\
\bottomrule
\end{tabularx}
\caption[Defensive-Offense layer architecture]{Layer architecture.
Core $\to$ Training direction is \textbf{prohibited}; Training $\to$
Core is permitted only via the Carryback procedure.}
\label{tab:def-offense-arch}
\end{table}

\subsection{Training v1.8 (Specification)}
\label{sec:training-v18}

\paragraph{Purpose.}
Avoid ruin in adversarial / pressure scenarios; analyse pattern
structures; generate defensive responses. \emph{No execution.}

\subsubsection{5.1 Input}
\label{sec:training-v18-input}

\begin{itemize}
\item ``Structure'' (e.g., scenario type, pattern ID).
\item Constraint (load, volatility, irreversibility flags).
\end{itemize}

\subsubsection{5.2 Output (Training)}
\label{sec:training-v18-output}

\begin{itemize}
\item Structure analysis (pattern recognition).
\item Defensive-Offense (defence design).
\item Prohibited-condition analysis (SOP-level conditions).
\item Conditions analysis (when shutdown / hold should be
considered).
\end{itemize}

\subsection{Carryback Specification}
\label{sec:carryback-spec}

Training outputs must traverse the Carryback procedure (4 steps)
before influencing operations. Direct/script/operation transfer is
\textbf{prohibited}.

\begin{enumerate}
\item \textbf{Tripwire}: ruin-related signal recognition (e.g., a
specific input pattern matches a Defensive-Offense category).
\item \textbf{Gate}: Type ZERO judgement gate (e.g., ``does this
trigger a mode transition?'').
\item \textbf{PP Template}: defensive design template recorded as a
Passive Pattern artefact.
\item \textbf{SOP Patch}: prohibited-condition / condition update at
the SOP layer (permitted-action set update, never direct execution).
\end{enumerate}

\paragraph{Effect.}
Training findings flow into the operational layer only as
\emph{constraints} (extending the prohibited set or refining the
permitted set), never as \emph{actions}.

\subsection{Prohibited (Hard Rules)}
\label{sec:def-offense-hard-prohibitions}

\begin{itemize}
\item Forwarding training output as operational action.
\item Forwarding ``persuasion patterns'' / scripted manipulation.
\item Direct training-to-operation use without Carryback.
\item ``It worked in training, so apply it'' --- this triggers
SHUTDOWN.
\end{itemize}

\subsection{v1.7 Notes (Historical, Two Items)}
\label{sec:def-offense-v17-notes}

\begin{enumerate}
\item \textbf{Defensive-Offense}: PP and Type ZERO must remain the
sole ``defensive'' interface; expansion is forbidden.
\item \textbf{Training}: v1.7 introduced ``Adversarial role-play''
with two clauses --- (a) it must not modify operational state, and
(b) it must not bypass the Carryback procedure. v1.8 formalises
Carryback (Tripwire / Gate / PP Template / SOP Patch).
\end{enumerate}

\paragraph{Signature.}
NRMO / Type ZERO (Revised Spec Draft) --- generated for operational
clarity. ``Design'' boundary is explicit: NRMO Core handles ruin
constraints; Type ZERO handles execution-action governance-layer
definition; SOP / TTM / HST-N handle their respective sub-domains.

\subsection{Positioning (Type ZERO Within NRMO)}
\label{sec:def-offense-positioning}

\paragraph{Action flow.}

\begin{lstlisting}[basicstyle=\ttfamily\small,frame=single]
[ Observation ]
       |
       v
   Action proposal A
       |
       v
[ Type ZERO ] -- ex-ante --> A_allowed definition
       |
       v
[ NRMO Core ] -- A_allowed evaluated against
                 ruin-constraint --> YES / NO / HOLD
       |
       v
   Selection / Execution
\end{lstlisting}

\begin{itemize}
\item NRMO Core: evaluates against the ruin constraint (time-path
survival).
\item Type ZERO: defines the action category (permitted set), based
on observed state.
\end{itemize}

\subsection{Design Invariants}
\label{sec:def-offense-design-invariants}

\begin{description}
\item[Ex-ante Governance] Action restriction occurs \emph{before}
NRMO Core is invoked, not after.
\item[Separation] Evaluation (NRMO) is structurally separated from
operation (Type ZERO) and training (SOP / TTM).
\item[Minimalism] No optimisation, no ranking, only audit-trail.
\item[Fail-Safe] Default is restriction, not permission.
\item[Non-Persuasion] No purpose-driven action shaping; no
narrative.
\end{description}

\subsection{Modes Summary}
\label{sec:def-offense-modes}

(Modes are defined formally in
Section~\ref{sec:type-zero-modes-c2}.)

\begin{itemize}
\item \textbf{Core} --- normal operation.
\item \textbf{Venture} --- exploration, bounded.
\item \textbf{Mission} --- goal-constrained, conditional commitment.
\item \textbf{Shutdown} --- defensive halt.
\end{itemize}

\subsection{I/O Specification}
\label{sec:def-offense-io}

\paragraph{Input (observation).}

\begin{lstlisting}[basicstyle=\ttfamily\small,frame=single]
Context:        (investment / work / relationship / etc.)
Buffer:         (capital / time / attention / etc.)
Observability:  observation strength
Irreversibility-Flag: irreversibility-imminent flag
Load:           load / fatigue
\end{lstlisting}

\paragraph{Output.}

\begin{lstlisting}[basicstyle=\ttfamily\small,frame=single]
Mode:        m in {Core, Venture, Mission, Shutdown}
A_allowed(m): permitted action category
Guardrails:  audit, conditions, restrictions
\end{lstlisting}

\subsection{Transition Rules}
\label{sec:def-offense-transition-rules}

\paragraph{Pseudocode (mode selection).}

\begin{lstlisting}[basicstyle=\ttfamily\small,frame=single]
if   shutdown_trigger():       mode = Shutdown
elif mission_armed():          mode = Mission
elif venture_conditions():     mode = Venture
else:                          mode = Core
\end{lstlisting}

\paragraph{Triggers.}

\begin{description}
\item[Shutdown-trigger]
\begin{itemize}
\item Buffer depleted.
\item Irreversibility imminent.
\item Observation state degraded.
\item Load saturated.
\end{itemize}
\item[Venture-conditions]
\begin{itemize}
\item Buffer sufficient.
\item Observability adequate.
\item Loss cap defined.
\end{itemize}
\item[Mission-armed]
\begin{itemize}
\item Mission condition (deadline, success criterion) defined.
\item Audit conditions met.
\item Irreversibility commitment justified.
\end{itemize}
\end{description}

\subsection{Hard Prohibitions}
\label{sec:def-offense-hard-rules}

\begin{itemize}
\item Direct execution by Type ZERO (no actuator authority).
\item In-play modification (Type ZERO does not modify itself during
operation).
\item Training requirement modification (training cannot modify SOP
/ TTM / specifications).
\end{itemize}

\subsection{Audit Items (Required Log Fields)}
\label{sec:def-offense-audit}

\begin{itemize}
\item Definition: 4 input variables.
\item $A_{\mathrm{allowed}}$ category definition (info / analysis /
proposal / etc.).
\item Shutdown-trigger occurrences (3--5 categories).
\item Mission / venture conditions and audits (definition /
deadline / criteria).
\end{itemize}

\paragraph{Version.}
Type ZERO Spec v1.0 (No-Psych) --- governance-layer definition only.
The version number is intentional: ``No-Psych'' indicates that the
Type ZERO governance layer does not perform psychological inference
(inference is reserved for HST-N).
